% TOP_PROBLEM: {"id": 3371, "title": "Is Skewes' number e^e^e^79 an integer?", "category_id": 1, "category": {"id": 1, "name": "number_theory", "display_name": "Number Theory", "description": "Properties of integers, prime numbers, Diophantine equations.", "slug": "number-theory", "order_index": 1, "created_at": "2026-07-31T15:26:25.670Z"}, "status": "open", "problem_number": "OPG-37366", "set_id": 12}
% FIRST_POSTED: 2026-10-01
% ENTRY_KIND: statement_only
% UnsolvedMath ID 3371; source catalogue code OPG-37366
% !TeX program = lualatex
% Definition and sources only; this file is not an LLM solution attempt.
\documentclass[11pt,a4paper]{article}
\usepackage[margin=25mm]{geometry}
\usepackage{fontspec}
\setmainfont{lmroman10-regular.otf}[BoldFont=lmroman10-bold.otf,
 ItalicFont=lmroman10-italic.otf,BoldItalicFont=lmroman10-bolditalic.otf]
\usepackage{amsmath,amssymb,mathtools}
\usepackage{unicode-math}
\setmathfont{latinmodern-math.otf}
\AtBeginDocument{\let\setminus\smallsetminus}
\usepackage{xurl,hyperref}
\hypersetup{colorlinks=true,urlcolor=blue}
\setcounter{secnumdepth}{0}
\setlength{\parindent}{0pt}
\setlength{\parskip}{5pt}
\setlength{\emergencystretch}{3em}
\begin{document}
\section*{Is Skewes' number $e^{e^{e^{79}}}$ an integer?}\label{3371}
{\small UnsolvedMath ID 3371 · OPG-37366 · Number Theory · OpenGarden}
\subsection{Definitions and mathematical statement}
Use the real exponential function $\exp:\mathbb R\to\mathbb R_{>0}$,
and define, successively,
\[
 a=\exp(79),\qquad b=\exp(a),\qquad X=\exp(b).
\]
Thus $X=e^{e^{e^{79}}}$ with exponentiation associated to the right.
The specific assertion in the source is $X\notin\mathbb Z$. An
equivalent formulation is that there exists an integer $m$ such that
$m<X<m+1$. Since $X>0$, only nonnegative integers are relevant when
asking whether it can equal an integer.

The title's reference to Skewes' number serves to identify this exact
exponential expression. The target is not the location of a sign
change in a prime-counting function, and changing the innermost
constant changes the number under discussion. Establishing that $X$
is irrational or transcendental would imply the desired conclusion,
but neither stronger assertion is demanded. A proof that an
intermediate quantity $a$ or $b$ is nonintegral does not by itself
settle nonintegrality after exponentiation. The problem asks for an
unconditional exact conclusion; numerical magnitude estimates, or
a consequence conditional on a further transcendence conjecture,
do not determine the integer-membership question.
\subsection{Short English statement}
Is the precisely specified real number exp(exp(exp(79))) nonintegral?
\subsection{Sources}
\begin{itemize}
\item[S1] ULAM.AI and the UnsolvedMath Contributors, supplied problems.json, record 3371 (OPG-37366), OpenGarden collection. Catalogue identity and source transcription; CC BY 4.0.
\url{https://huggingface.co/datasets/ulamai/UnsolvedMath}.
\item[S2] Open Problem Garden, Is Skewes' number e\textasciicircum{}e\textasciicircum{}e\textasciicircum{}79 an integer?. Original problem and discussion; the mathematical formulation below is expanded from the supplied source record.
\url{http://www.openproblemgarden.org/op/is_skewes_number_e_e_e_79_an_integer}.
\end{itemize}
\end{document}
